\section{Related Work}
\label{sup:related}

\runin{Testing theory.} Baarda's data snooping and reliability theory and
its extension to the detection--identification--adaptation (DIA) procedure test a linear Gaussian
model with nuisance parameters for a misspecified observation and give a minimum detectable bias
\cite{stat_teunissen_dia}. In satellite navigation, solution separation, which compares the full
solution with a leave-one-out solution, equals the residual-based test \cite{flt_raim_joerger}.

\runin{Message passing.} Expectation propagation approximates a posterior by a product of site
messages, and its cavity step removes one site's message before re-fitting it by moment matching
\cite{mp_minka_ep}; for Gaussian messages the removal is an exact subtraction in information form.
Gaussian belief propagation passes such messages on factor graphs and has been proposed as a
computational basis for spatial AI \cite{mp_davison_gbp}. Our fuse step is one round of Gaussian messages
from views to surfels, and our leave-view-out is the cavity step alone, without the moment-matching
iterations of full expectation propagation.

\runin{Our gate in these terms.} The gate is an instance of Baarda's testing theory; its observations are pixels,
its nuisance is the per-view scale and offset, and its leave-one-out solution is a rendered surface.

\runin{Monocular depth.} Metric depth models \cite{est_metric3dv2,est_unidepth,est_depthpro} and
affine-invariant ones \cite{est_midas,est_dav2,est_marigold,est_moge} are used as priors for
reconstruction, usually after a per-view scale and shift alignment \cite{rec_monosdf}. Their
failures on reflective, transparent and textureless surfaces are known; whether a model repeats the
same error across views is what H2 measures.

\runin{Surfels and Gaussian splatting.} Point-based fusion maintains surfels with a normal, a radius
and a confidence and fuses depth by confidence-weighted averaging with projective association
\cite{fus_keller_surfels}. 3D Gaussian splatting renders anisotropic Gaussians with a differentiable rasteriser
\cite{gs_3dgs}; 2DGS uses planar Gaussians for accurate surfaces \cite{gs_2dgs}; gsplat is an
open-source rasteriser \cite{gs_gsplat}; DN-Splatter adds depth and normal priors \cite{gs_dnsplatter};
3DGS-MCMC replaces densification by relocation \cite{gs_mcmc}. Our surfels are 2DGS surfels with a
separate uncertainty block, fused by the rasteriser's own Jacobians instead of by association.

\runin{Robust reconstruction and distractors.} Robust losses \cite{rob_blackrangarajan,rob_barron}
and distractor-aware methods such as RobustNeRF \cite{gs_robustnerf}, SpotLessSplats
\cite{gs_spotless} and RobustSplat \cite{gs_robustsplat} down-weight pixels or regions whose
photometric residual looks like an outlier. They act per pixel or per region, mostly from colour,
and they state no false-alarm rate. Multi-view depth consistency filters \cite{fus_merrell,fus_colmap_mvs}
test pixels against other views' depth maps by correspondence.

\runin{Uncertainty in radiance fields.} FisherRF computes Fisher information of radiance-field
parameters for active view selection and uncertainty quantification \cite{unc_fisherrf}; Bayes' Rays
estimates a post-hoc spatial uncertainty of a trained NeRF with a Laplace approximation over a
perturbation field and uses it, among other things, to remove artefacts \cite{unc_bayesrays}. Neither,
to our knowledge, keeps a per-primitive uncertainty that renders into depth variance and is used,
with exact leave-view-out, to test whole views at a stated false-alarm rate. H6 uses both as
alternative uncertainties inside the same test. \TODO{Final literature pass over uncertainty-aware
Gaussian splatting and depth-prior filtering before the preprint.}
